% GENERATED FILE. Edit canonical lessons, then run npm run generate:books.
% canonical-source-hash: fnv1a64:2921525d0ce2b056
% canonical-lessons: KA-C45-grandfather, KA-C45-grandmother, KA-C45-husband, KA-C45-wife

\chapter{The Older Generation}
\label{ch:ka-older-generation}
\hlchaptermodality{45}

\begin{chapteropening}
\textbf{By the end of this chapter:} \emph{I can name a grandfather, a grandmother, a husband and a wife in Kannada.}

Complete KA-C45-wife and say all four words in order.
\end{chapteropening}

\section[ajja]{\kn{ಅಜ್ಜ} (ajja) — grandfather}

\label{lesson:KA-C45-grandfather}



\begin{quote}
Before the new one: what did \emph{bare} mean?
\end{quote}

\subsection*{You'll want to know: \kn{ಅಜ್ಜ}}
\textbf{\kn{ಅಜ್ಜ}} (\emph{ajja}) — \textquotedblleft{}grandfather\textquotedblright{}.

A native Dravidian word of the short, doubled shape that small children's mouths reach for first, and it carries respect well past the family: an unrelated old man can be \emph{ajja} too. The \emph{-a} closing it is the masculine ending.

\begin{grammarlens}[title={What you've built}]
The first of four words for the generation above you and the one beside you. Three follow, and each reuses the ones before.
\end{grammarlens}

\subsection*{Your turn}
Say these aloud:

\begin{itemize}
\raggedright
  \item \emph{ajja}
  \item it once more, slowly
  \item it after \emph{bare}, so the two sit together
\end{itemize}

\begin{tcolorbox}[breakable,colback=teal!4,colframe=teal!35!black,title={Before you move on}]
What does \kn{ಅಜ್ಜ} mean? (\textquotedblleft{}grandfather\textquotedblright{}.) Say it once more.
\end{tcolorbox}
\section[ajji]{\kn{ಅಜ್ಜಿ} (ajji) — grandmother}

\label{lesson:KA-C45-grandmother}



\begin{quote}
Before the new one: what did \emph{ajja} mean?
\end{quote}

\subsection*{You'll want to know: \kn{ಅಜ್ಜಿ}}
\textbf{\kn{ಅಜ್ಜಿ}} (\emph{ajji}) — \textquotedblleft{}grandmother\textquotedblright{}.

The same stem as the last word, with the ending swapped: \emph{-i} where \emph{ajja} had \emph{-a}. The son-and-daughter pair built \kn{ಮಗಳು} out of \kn{ಮಗ} by that same trick, one stem carrying two endings, and this pair is tidier still — only the final vowel moves.

\begin{grammarlens}[title={What you've built}]
Two of the four, and a pattern you have now seen twice in this book.
\end{grammarlens}

\subsection*{Your turn}
Say these aloud:

\begin{itemize}
\raggedright
  \item \emph{ajji}
  \item it once more, slowly
  \item it after \emph{ajja}, and hear only the last vowel move
\end{itemize}

\begin{tcolorbox}[breakable,colback=teal!4,colframe=teal!35!black,title={Before you move on}]
What does \kn{ಅಜ್ಜಿ} mean? (\textquotedblleft{}grandmother\textquotedblright{}.) Say it once more.
\end{tcolorbox}
\section[gaṇḍa]{\kn{ಗಂಡ} (gaṇḍa) — husband}

\label{lesson:KA-C45-husband}



\begin{quote}
Before the new one: what did \emph{ajji} mean?
\end{quote}

\subsection*{You'll want to know: \kn{ಗಂಡ}}
\textbf{\kn{ಗಂಡ}} (\emph{gaṇḍa}) — \textquotedblleft{}husband\textquotedblright{}.

The everyday word, and a native one. Its own stem also gives \kn{ಗಂಡು} (\emph{gaṇḍu}), 'male', so what it names is a man of the house rather than anything about the wedding. Formal Kannada keeps the Sanskrit \emph{pati} alongside it for writing and ceremony.

\begin{grammarlens}[title={What you've built}]
Three of the four. The last one is built the same way, from the other half of the pair.
\end{grammarlens}

\subsection*{Your turn}
Say these aloud:

\begin{itemize}
\raggedright
  \item \emph{gaṇḍa}
  \item it once more, slowly
  \item it after \emph{ajji}, so the two sit together
\end{itemize}

\begin{tcolorbox}[breakable,colback=teal!4,colframe=teal!35!black,title={Before you move on}]
What does \kn{ಗಂಡ} mean? (\textquotedblleft{}husband\textquotedblright{}.) Say it once more.
\end{tcolorbox}
\section[heṇḍati]{\kn{ಹೆಂಡತಿ} (heṇḍati) — wife}

\label{lesson:KA-C45-wife}



\begin{quote}
Before the new one: what did \emph{gaṇḍa} mean?
\end{quote}

\subsection*{You'll want to know: \kn{ಹೆಂಡತಿ}}
\textbf{\kn{ಹೆಂಡತಿ}} (\emph{heṇḍati}) — \textquotedblleft{}wife\textquotedblright{}.

Built on \kn{ಹೆಣ್ಣು} (\emph{heṇṇu}), 'female', with a \emph{-ti} ending. And \emph{heṇṇu} carries the sound change you already met in \kn{ಹೋಗು}: where its sisters keep an old \emph{p-}, Kannada has \emph{h-}. So Tamil \emph{peṇ}, 'woman', and this \emph{heṇ-} are one word wearing two faces.

\begin{grammarlens}[title={What you've built}]
That closes the run: grandfather, grandmother, husband, wife.
\end{grammarlens}

\subsection*{Your turn}
Say these aloud:

\begin{itemize}
\raggedright
  \item \emph{heṇḍati}
  \item it once more, slowly
  \item it after \emph{gaṇḍa}, so the two sit together
\end{itemize}

\begin{tcolorbox}[breakable,colback=teal!4,colframe=teal!35!black,title={Before you move on}]
What does \kn{ಹೆಂಡತಿ} mean? (\textquotedblleft{}wife\textquotedblright{}.) Say it once more.
\end{tcolorbox}
